\section{Method}
\label{sec:method}

\subsection{Parametric pre-compensation}
\label{sec:parametric}

The interpretable pre-compensation is a filter chain implemented with FFmpeg~\cite{ffmpeg}:
spatio-temporal denoising (\texttt{hqdn3d}), contrast and saturation (\texttt{eq}), Lanczos
upscaling, an unsharp mask on luma, and temporal luma grain. \Cref{tab:space} lists the seven
parameters. All are continuous; the upscale factor is snapped to $\{1, 1.5, 2\}$ and the unsharp
radius to $\{1,2,3\}$ (kernel size $3,5,7$), so the same space serves every sampler. The upscale
factor is the tier component $s$ of \cref{def:decomposition}; the other six form the filter
component $\varphi$. The upload is encoded near-losslessly (x264, CRF 10), so the uploader's own
compression is not a confounder.

\begin{table}[t]
  \centering
  \caption{Parameters of the interpretable pre-compensation. The identity value makes a step a
  no-op; trial 0 of every search is the identity (the naive upload).}
  \label{tab:space}
  \begin{tabular}{@{}llll@{}}
    \toprule
    Parameter & Range & Identity & Role \\
    \midrule
    denoise        & $[0, 6]$       & 0   & \texttt{hqdn3d} strength \\
    contrast       & $[0.90, 1.15]$ & 1   & \texttt{eq} contrast \\
    saturation     & $[0.85, 1.25]$ & 1   & \texttt{eq} saturation \\
    upscale        & $\{1, 1.5, 2\}$ & 1  & upload resolution factor (tier) \\
    sharpen        & $[0, 2]$       & 0   & unsharp amount on luma \\
    sharpen\_radius & $\{1, 2, 3\}$ & 1   & unsharp kernel $3\times3$, $5\times5$, $7\times7$ \\
    grain          & $[0, 15]$      & 0   & temporal luma noise (dither) \\
    \bottomrule
  \end{tabular}
\end{table}

\paragraph{Search.}
We compare Gaussian-process Bayesian optimisation (GP-BO, 8 start-up trials), multivariate
TPE (8 start-up trials), CMA-ES (population 8) and random search, all from
Optuna~\cite{akiba2019optuna}, with a budget of 40 trials. A trial renders the upload, sends it
through the platform and scores the playback on every clip of the training set; the objective is
the mean over clips. Evaluations are cached by $(\text{clip},\theta)$, so studies, transfers and
re-scoring share work.

\subsection{Learned pre-compensation through an emulator}
\label{sec:learned}

\paragraph{Networks.}
Both the emulator $\Em$ and the pre-compensation $\Gt$ are residual U-Nets~\cite{ronneberger2015unet}
of the form $y = x + r\tanh(h(x))$, with the last convolution initialised to zero so that each
network starts as the identity. The emulator has width 32, three levels and $r=1$. The
pre-compensation has width 24, two levels and a hard bound $r=\delta=0.06$ on its per-pixel change
(on a $[0,1]$ scale). Both operate on RGB frames; the upload tier is fixed at $\times 2$ for the
learned route, since the tier is settled by the parametric study.

\paragraph{Training.}
In round 0 we send, for each training clip, the identity and five random parametric
pre-compensations (at the source resolution) through the black box, and keep every fourth frame of
each (upload, playback) pair. The emulator is trained with an $\ell_1$ loss on aligned
$256\times256$ crops (AdamW, 3000 steps). The pre-compensation is then trained through the frozen
emulator with
\begin{equation}
  \mathcal{L}(\theta) = \mathrm{LPIPS}\big(\Em(\Gt(x)), x\big)
  + \big\|\Em(\Gt(x)) - x\big\|_1
  + \lambda\,\big\|\Gt(x) - x\big\|_1,
  \label{eq:gloss}
\end{equation}
with LPIPS (AlexNet)~\cite{zhang2018unreasonable}, $\lambda=2$ (the trust region of
\cref{eq:trust}) and 1500 Adam steps. In P1 the perceptual term can also be
DISTS~\cite{ding2022image} or the sum of both. After each round, $\Gt$'s outputs are sent through the
real black box; their real playback is added to the emulator's data for the next round (domain
adaptation), and the emulator's PSNR against the real playback is logged as the measured
sim-to-real gap.

\paragraph{Real-anchored estimator.}
\label{sec:anchored}
To test whether $\Gt$ fails because it exploits the emulator (\cref{sec:emulator-theory}), we also
train it with a forward value anchored on the black box. Training is split into $A$ phases. At the
start of phase $k$, the current $G_k$ is applied to every training clip and sent once through the
black box, giving $y^{\text{real}}_k=\fp(G_k(x))$, and the offset
$\Delta_k = y^{\text{real}}_k - \Em(G_k(x))$ is stored. During the phase the loss of \cref{eq:gloss}
is evaluated on
\begin{equation}
  \hat y = \Em(\Gt(x)) + \sg(\Delta_k),
  \label{eq:anchored}
\end{equation}
which equals the black box's output at the anchor, adds the emulator's predicted change away from
it, and back-propagates the emulator's gradient. This is the straight-through estimator
$y^{\text{real}} + \Em(x) - \sg(\Em(x))$~\cite{bengio2013estimating,tian2021self,pesnel2026scaled} with the real output refreshed once per
phase; a per-step refresh would cost one upload per clip and step. Each anchor costs one black-box
call per training clip.

\paragraph{Pre-compensation on top of a tuned configuration.}
To ask whether $\Gt$ can add to a tuned interpretable configuration $\theta^\ast$ rather than compete
with it, we move $\theta^\ast$ into the channel: the upload is $\theta^\ast(\Gt(x))$, a new emulator
of $\fp\circ\theta^\ast$ is trained, and $\Gt$ is trained with the anchored estimator. At
$\Gt=\mathrm{id}$ the result is exactly $\theta^\ast$'s, so any change is what $\Gt$ adds.

\subsection{Fixed recipes and upload bitrate}
\label{sec:recipes}
As baselines that need no tuning, and no uploads to tune, we evaluate fixed recipes: $\times2$ with
FFmpeg's default unsharp ($5\times5$, amount 1.0), $\times2$ with a mild unsharp ($5\times5$, 0.5),
and both the tier-only upload and the mild recipe encoded at a fixed 16\,Mb/s (x264, rate-capped)
instead of near-losslessly.

\subsection{Metrics}
\label{sec:metrics}

All metrics compare the playback, resampled bicubically to the source resolution, with the source.

\paragraph{Objective.}
The default objective is
\begin{equation}
  \text{composite} = \text{VMAF}_{\text{NEG}} - 40\cdot\text{LPIPS} - 2\cdot\text{CAMBI},
\end{equation}
combining the no-enhancement-gain variant of VMAF~\cite{li2016toward}, LPIPS and the CAMBI banding
index~\cite{tandon2021cambi}. VMAF NEG is used because plain VMAF rewards sharpening, which an
optimiser would exploit. The weights were set by hand to put the terms on a comparable scale.
\todo{justify the weights or report sensitivity to them.}

\paragraph{Judges.}
Three metrics are never optimised and serve as judges (\cref{def:confirmed}):
SSIMULACRA~2~\cite{sneyers2022ssimulacra2} (every third frame), DISTS~\cite{ding2022image} (every
third frame) and luma XPSNR~\cite{helmrich2020xpsnr} (FFmpeg \texttt{xpsnr} filter). PSNR and SSIM
are recorded for reference.
